\begin{frame}{확인 문제 (9.1)}
  \small
  \begin{enumerate}
    \item \kb{XOR의 부등식 모순.} $w_2+b>0$, $w_1+b>0$, $b\le 0$,
      $w_1+w_2+b\le 0$ — 첫 두 식을 더해 셋째 식과 모순을 만들어,
      이 네 식을 모두 만족하는 $(w_1,w_2,b)$가 \textbf{존재하지
      않음}을 보여라.
    \item \kb{$\delta_1$의 재사용.} $\delta_1 = (W_2\delta_2)\odot
      \sigma'(z_1)$에서 $\delta_2$가 출력층에서 이미 계산돼 \textbf{재사용}
      된다는 것이 역전파가 ``동적 프로그래밍''인 이유. 10층 망에서
      모든 $\delta_l$을 처음부터 다시 계산하면 몇 번의 미분이 필요한가,
      재사용하면 몇 번인가?
    \item \kb{교차 엔트로피의 약분.} 시그모이드 + 교차 엔트로피면
      $\delta_2 = a_2 - y$로 미분항이 약분. MSE 때의 $\delta_2$와
      비교하고, 그 차이가 ``포화된 예측에서도 그래디언트가 살아있는''
      것과 어떻게 연결되는지 한 문장으로.
  \end{enumerate}
\end{frame}
